\section{Estimation and screening theory}\label{sec:estimation}

\subsection{Empirical-rank local Hill scores}

The estimator follows the envelope construction coordinate by coordinate.
Let $R_{ij}$ be the rank of $X_{ij}$ among $X_{1j},\ldots,X_{nj}$ and
$\widehat U_{ij}=R_{ij}/(n+1)$. For $u\in\I$, form the rank window
\[
 \mathcal W_j(u)=\{i:|\widehat U_{ij}-u|\le h\},\qquad
 M_j(u)=|\mathcal W_j(u)|,\qquad
 k_j(u)=\lfloor\alpha M_j(u)\rfloor,
\]
and, writing $Y_{j,u,(1)}\le\cdots\le Y_{j,u,(M_j(u))}$ for the local order
statistics, the local Hill estimate
\begin{equation}\label{eq:hill}
 \widehat\xi_j(u)=\frac1{k_j(u)}\sum_{r=1}^{k_j(u)}
 \log\frac{Y_{j,u,(M_j(u)-r+1)}}{Y_{j,u,(M_j(u)-k_j(u))}},
\end{equation}
defined when $\alpha M_j(u)>1$. This is the local-fraction Hill statistic of
\citet[Definition~5]{gardesPodgorny2025} evaluated on one-coordinate rank
windows. The score averages the profile over the deterministic grid
$G=\{r/(n+1):r=1,\ldots,n\}\cap\I$,
\begin{equation}\label{eq:emp-score}
 \widehat\Psi_j=\frac1{|G_j^{\rm val}|}\sum_{u\in G_j^{\rm val}}\widehat\xi_j(u),
 \qquad
 G_j^{\rm val}=\{u\in G:\alpha M_j(u)>1\},
\end{equation}
Because the window counts $M_j(u)$ depend on the ranks only through
the grid position, $G_j^{\rm val}$ is the same set for every $j$, and
under the conditions of Theorem~\ref{thm:score} every grid point is
usable, so $G_j^{\rm val}=G$; the convention $\widehat\Psi_j=+\infty$
when no grid point is usable is an implementation safeguard that is
provably never triggered when $K_n\ge2$. The screened set
$\widehat\A_d$ collects the $d$ smallest scores, ties broken by a
seeded auxiliary randomization.

The two localization parameters play different roles: $h$ fixes how many
observations resemble the target rank position, roughly $2nh$ of them, and
$\alpha$ fixes how many of their largest responses enter the Hill statistic,
roughly $2n\alpha h$. Rank windows make these counts deterministic up to
boundary effects, which is what permits exact finite-sample union bounds
over a diverging number of coordinates.

\subsection{Assumptions}

Ranks displace the unobserved transforms $U_{ij}$, windows displace the
conditioning point, and the Hill statistic sees only the upper
$\alpha$-fraction of a projected distribution. The assumptions bound these
three effects uniformly over the coordinates. Throughout this section
every active coordinate is assumed detectable, $\Delta_{\min}>0$ for
every $n$, so that the conditions below are well posed. Let
$Q_j(\cdot,u)$ denote the generalized upper conditional quantile of $Y$
given $U_j=u$, $Q_j(s,u)=F_{Y|j}^{-1}(1-s^{-1}\mid u)$ with
$F^{-1}(t)=\inf\{y:F(y)\ge t\}$, and let $\tau_n=(pn)^{-2}$.

\begin{itemize}
\item[(E1)] \emph{Quantile representation and regularity of the projected
index.} For every $j\le p$ and every $u\in[0,1]$ the conditional
distribution function $F_{Y|j}(\cdot\mid u)$ is continuous, so that
$V_{ij}=1-F_{Y|j}(Y_i\mid U_{ij})$ is conditionally standard uniform and
$Y_i=Q_j(1/V_{ij},U_{ij})$ almost surely. Moreover the projected indices
are equicontinuous on $\I$, uniformly over coordinates and over $n$,
\[
 \omega_\xi(r)=\max_{j\le p}\;
 \sup_{u,v\in\I,\,|u-v|\le r}|\xi_j(u)-\xi_j(v)|
 \longrightarrow0\qquad(r\to0),
\]
and satisfy the grid-scale regularity condition
$\omega_\xi(1/n)=o(\Delta_{\min})$.
\item[(E2)] \emph{Uniform tail regularity.} Writing
$Q_j(s,u)=s^{\xi_j(u)}\ell_j(s,u)$ with $\ell_j(\cdot,u)$ slowly varying,
\[
 \max_{j\le p}\,\sup_{u\in\I}\,
 \sup_{s\ge(3\alpha)^{-1},\,t>1}
 \frac{|\log\{\ell_j(ts,u)/\ell_j(s,u)\}|}{\log t}
 \;=\;o(\Delta_{\min}).
\]
\item[(E3)] \emph{Uniform local quantile regularity.} There exists a
constant $C_0>1$ such that
\[
 \max_{j\le p}\,
 \sup_{u\in\I}\;
 \sup_{v\in[0,1],\,|u-v|\le C_0h}\;
 \sup_{w\in[\tau_n,1-\tau_n]}
 \left|\log\frac{Q_j(w^{-1},v)}{Q_j(w^{-1},u)}\right|
 \;=\;o(\Delta_{\min}).
\]
The theorems use only this one $C_0$; under \eqref{eq:rate-cond} any
$C_0>1$ suffices.
\end{itemize}


The condition (E2) is the analogue of condition (H.4) in
\citet{gardesPodgorny2025}, normalized by $\log t$ and required only beyond
the working tail level $(3\alpha)^{-1}$. (E3)
adapts their condition (H.5): the truncation $\tau_n=(pn)^{-2}$ makes the
probability that any of the $pn$ conditional PIT variables leaves
$[\tau_n,1-\tau_n]$ vanish, which is the high-dimensional replacement for
their fixed-dimension truncation. These conditions are stated on the projected quantiles, not derived from
(C1)--(C2): projection can generate genuinely new slow variation. If
$W\sim U(0,1)$ and $\Pp(Y>y\mid W)=y^{-1/(1+W)}$ exactly, then
$\Pp(Y>y)\sim4y^{-1/2}/\log y$, so even exact Pareto conditionals produce
a logarithmic factor after projection, and the left side of (E2) is then
of exact order $1/\log(1/\alpha)$. This logarithmic effect also suggests that the usual polynomially shrinking signal assumptions from the SIS literature may be too optimistic in the present EVT setting. In particular, if $\alpha=n^{-a}$, the preceding example produces a projected-tail bias of order $1/\log(1/\alpha)\asymp1/\log n$. Condition (E2) then requires this bias to be negligible relative to $\Delta_{\min}$, which is incompatible with a gap of polynomial order $\Delta_{\min}\asymp n^{-\kappa}$ for any $\kappa>0$, since $n^{-\kappa}=o(1/\log n)$. We therefore expect that, for realistic extreme-value models in which projection naturally generates logarithmic slow variation, admissible signal gaps may have to vanish substantially more slowly than the polynomial rates commonly imposed in classical SIS theory.

\subsection{Main results}

\begin{theorem}[Score concentration]\label{thm:score}
Assume (C1)--(C2), (S), (E1)--(E3), $\alpha\le1/3$,
$h\le1-\varepsilon$, and
\begin{equation}\label{eq:rate-cond}
 \log(pn)=o(nh^2),
 \qquad
 \log(pn)=o(n\alpha h).
\end{equation}
Then there are deterministic sequences
$a_n=O(\sqrt{\log(pn)/(n\alpha)})$ and $\beta_n=o(\Delta_{\min})$ with
\[
 \Pp\Bigl(\max_{j\le p}\,|\widehat\Psi_j-\Psi_j|\le a_n+\beta_n\Bigr)
 \longrightarrow1;
\]
in shorthand,
$\max_{j\le p}|\widehat\Psi_j-\Psi_j|
=O_{\Pp}(\sqrt{\log(pn)/(n\alpha)})+o(\Delta_{\min})$.
\end{theorem}

The bandwidth does not appear in the stochastic rate. It is worth being
precise about why, because the natural guess is that it should. Each
local Hill statistic is built from about $n\alpha h$ spacings and does
fluctuate at scale $(n\alpha h)^{-1/2}$; but $\widehat\Psi_j$ is not a
local statistic, it is an \emph{average} of $O(n)$ of them over the rank
grid, and the fluctuation of an average is not the largest of the
fluctuations.

The $o(\Delta_{\min})$ term collects the tail bias (E2),
the spatial bias (E3) evaluated at $C_0h$ --- the first condition in
\eqref{eq:rate-cond} makes the rank displacement smaller than $(C_0-1)h$,
so windows never leave the range covered by (E3) --- and the quadrature
error $\omega_\xi(1/n)+O(1/n)$ from (E1).

\begin{theorem}[Sure screening]\label{thm:sis}
Assume the conditions of Theorem~\ref{thm:score} (including
$\Delta_{\min}>0$, i.e.\ every active coordinate detectable in the
sense of Proposition~\ref{prop:geometry}) and
\begin{equation}\label{eq:sis-cond}
 \log(pn)=o\bigl(n\alpha\,\Delta_{\min}^2\bigr).
\end{equation}
Then, for every $d\ge s$,
\[
 \Pp\bigl(\A\subseteq\widehat\A_d\bigr)\longrightarrow1,
\]
and with probability tending to one every active score is smaller than
every inactive score.
\end{theorem}


Condition \eqref{eq:sis-cond} quantifies the cost of dimension:
screening $p$ coordinates costs $\log p$ against the tail count
$n\alpha$, deflated by the squared gap. The bandwidth has left this
condition, but it has not left the theorem: it survives in
\eqref{eq:rate-cond}, where $\log(pn)=o(n\alpha h)$ is what makes every
local threshold lie inside the tail region, and $\log(pn)=o(nh^2)$
controls the rank displacement. The gain is therefore a
\emph{decoupling} rather than a removal.

Writing $\alpha=n^{-a}$, $h=n^{-b}/2$ and $\Delta_{\min}\asymp n^{-d}$,
the previous analysis required the single coupled restriction
$a+b+2d<1$. It is replaced by the two separate restrictions
\[
 a+b<1
 \qquad\text{and}\qquad
 a+2d<1,
\]
together with $b<1/2$ from the rank-displacement condition, and $b>d$
whenever the spatial bias is Lipschitz and must satisfy
$h_n=o(\Delta_{\min,n})$. For fixed gaps, $d=0$, the admissible region
does not enlarge, because $a+b<1$ still binds; the enlargement is strict
as soon as the gap shrinks. If instead $\log p\asymp n^\kappa$ grows
subexponentially, the three conditions read $\kappa<1-2b$,
$\kappa<1-a-b$ and $\kappa<1-a-2d$, in place of the former
$\kappa<1-a-b-2d$.
